\documentclass{article}

\usepackage[english]{babel}

\usepackage[letterpaper,top=2cm,bottom=2cm,left=3cm,right=3cm,marginparwidth=1.75cm]{geometry}

% Useful packages
\usepackage{amsmath}
\usepackage{graphicx}
\usepackage[colorlinks=true, allcolors=blue]{hyperref}
\usepackage{amsthm}
\usepackage{csquotes}
\usepackage{algorithm, algpseudocode, float}


\theoremstyle{definition}
\newtheorem{definition}{Definition}[section]


\title{COSC 473 Reading Responses}
\author{
    Max Moir \\
    \textit{University of Canterbury} \\ 
    \text{Christchurch, New Zealand} \\ 
    \texttt{maxwmoir@gmail.com}
}

% The proposal (PDF with 1-3 pages excluding references; single-column text) should identify the problem you wish to address and the motivation for doing so.
% It should also identify prior attempts to solve the problem (or a related problem) and analyse why these attempts were insufficient.
% The proposal should also provide an outline of the intended method for resolving the problem.
% Please also discuss risks that may affect your project (for example, alert-level changes or unavailability of data/tools your project depends on) and how you will mitigate the risks.

\begin{document}
\maketitle

\section{Introduction}

Reading response answers

\subsection{Week 4}
IPFS was designed to be a decentralized, content-addressed, peer-to-peer foundation for the future of the web.

\noindent
Based on the readings discuss the following:

\noindent
\emph{To what extent does the current reality of IPFS live up to its original idea?}

\noindent
Consider one way IPFS fulfills its vision and one way it falls short.

\subsection{Week 5}
In a centralised system, responsibility for knowing where data is concentrated, but Kademlia distributes it.
Kademlia is designed so that no one can control it. When is that a strength, and when might it be a problem?

Think about where else in life or society we rely on local knowledge and word-of-mouth to make things work without a central coordinator.

\subsection{Week 5}

What kinds of trust are present in Bitcoin, and how do they compare to the forms of trust in earlier systems or traditional finanical institutions?
How does Bitcoin establish agreement on a shared history of transactions and how does that compare to Lamport's way of reasoning about order in distributed systems?





\end{document}
